\section{Week 1: Protocol - Foundations of Representation}

This is the operational mandate for the first seven days. The objective is singular: establish a rhythm of deep work in the two core pillars of this year—Computer Systems and Mathematics. Distractions are eliminated. The 2-4 hour block is a non-negotiable commitment.

Adherence to the schedule is the only metric for success this week. The work must be done.

\noindent\rule{\textwidth}{0.4pt}

\subsubsection{\textbf{Monday: Computer Science - The Language of Memory}}
\begin{itemize}
  \item \textbf{Objective:} Master hexadecimal notation, the concept of "word size," and the critical distinction between big-endian and little-endian byte ordering.
  \item \textbf{Block (2 Hours):} Read CS:APP, pp. 29-41 (Sección 2.1.1 a 2.1.4).
  \item \textbf{Execution:}
\end{itemize}
\begin{enumerate}
  \item Escribe un programa en C que declare una variable \texttt{int x = 0x01234567;} y luego imprima sus bytes individuales para determinar el endianness de tu máquina.
  \item Convierte a mano los números decimales 15, 16, 255, y 256 a binario y hexadecimal.
  \item Explica en tus notas por qué el byte ordering es un problema fundamental en la programación de redes.

\end{enumerate}
\subsubsection{\textbf{Tuesday: Mathematics - The Geometry of Equations}}
\begin{itemize}
  \item \textbf{Objective:} To visualize linear equations not as abstract algebra, but as geometric objects: lines, planes, and their intersections.
  \item \textbf{Block (2-3 Hours):} Watch MIT 18.06 Linear Algebra, Lecture 1 ("The Geometry of Linear Equations") and Lecture 2 ("Elimination with Matrices").
  \item \textbf{Execution:}
\end{itemize}
\begin{enumerate}
  \item Take handwritten notes. The physical act of writing aids memory.
  \item For a 2x2 system, draw the row picture (intersecting lines) and the column picture (vector addition). Do this for a system with one solution, no solution, and infinite solutions.
  \item Perform Gaussian elimination on a 3x3 system by hand. Show your work.

\end{enumerate}
\subsubsection{\textbf{Wednesday: Computer Science - Bit-Level Manipulation}}
\begin{itemize}
  \item \textbf{Objective:} Entender cómo se representan las cadenas de caracteres y el código, y dominar las operaciones a nivel de bit en C.
  \item \textbf{Block (2 Hours):} Read CS:APP, pp. 46-56 (Sección 2.1.5 a 2.1.10).
  \item \textbf{Execution:}
\end{itemize}
\begin{enumerate}
  \item Escribe una función en C que use máscaras de bits (\texttt{\&}) y desplazamientos (\texttt{\textgreater{}\textgreater{}}) para extraer el byte más significativo de un entero de 32 bits.
  \item Implementa la operación XOR (\texttt{\textasciicircum{}}) usando solo las operaciones AND (\texttt{\&}) y NOT (\texttt{\textasciitilde{}}).
  \item Explica la diferencia fundamental entre el desplazamiento lógico a la derecha y el desplazamiento aritmético a la derecha.

\end{enumerate}
\subsubsection{\textbf{Thursday: Mathematics - Elimination and Inverses}}
\begin{itemize}
  \item \textbf{Objective:} To master the mechanics of matrix operations, the core machinery of linear algebra.
  \item \textbf{Block (2-3 Hours):} Watch MIT 18.06, Lecture 3 ("Multiplication and Inverse Matrices") and Lecture 4 ("Factorization into A = LU").
  \item \textbf{Execution:}
\end{itemize}
\begin{enumerate}
  \item Multiply two 3x3 matrices by hand.
  \item Find the inverse of a 2x2 matrix using the formula.
  \item Take a 3x3 matrix and perform LU decomposition on it by hand.

\end{enumerate}
\subsubsection{\textbf{Friday: The Forge - Armory and First Contact}}
\begin{itemize}
  \item \textbf{Objective:} To prepare your development environment and make first contact with your primary language. A warrior's weapon must be sharp and ready.
  \item \textbf{Block (2-3 Hours):}
  \item \textbf{Execution:}
\end{itemize}
\begin{enumerate}
  \item Install a C compiler (GCC via MinGW, macOS Command Line Tools, or \texttt{build-essential} on Linux).
  \item Install Rust via \texttt{rustup}.
  \item Write, compile, and run \texttt{hello, world} in C.
  \item Write, compile, and run \texttt{hello, world} in Rust using \texttt{rustc} directly.
  \item Create a new project using \texttt{cargo new hello\_cargo}, build it, and run it. Read the output to understand what Cargo is doing for you.

\end{enumerate}
\subsubsection{\textbf{Saturday: The Forge - Basic Training}}
\begin{itemize}
  \item \textbf{Objective:} To move from "hello, world" to programs that perform logic. Solidify basic syntax and concepts.
  \item \textbf{Block (3-4 Hours):} Work through Chapter 3 of \textit{The Rust Programming Language} ("Common Programming Concepts").
  \item \textbf{Execution:}
\end{itemize}
\begin{enumerate}
  \item For every concept (variables, data types, functions, control flow), write a small, separate program in a \texttt{practice} directory to test it.
  \item Build the three programs requested in the book:
\end{enumerate}
\begin{itemize}
  \item Temperature converter.
  \item Fibonacci number generator.
  \item "The Twelve Days of Christmas" lyrics generator.

\end{itemize}
\subsubsection{\textbf{Sunday: Review and Reload}}
\begin{itemize}
  \item \textbf{Objective:} To assess the week's progress, identify weak points, and prepare for the next operational cycle.
  \item \textbf{Block (2 Hours):}
  \item \textbf{Execution:}
\end{itemize}
\begin{enumerate}
  \item Read through every single note you took this week.
  \item Create a "Weak Points" list. Write down the 1-3 concepts that you feel least confident about.
  \item Create the \texttt{week\_02.md} file. The first task for next week will be to attack the items on your "Weak Points" list.
  \item Outline the new topics for Week 2 in Year1.md

\end{enumerate}
\noindent\rule{\textwidth}{0.4pt}

The plan is updated. The mission is unchanged.
